\documentclass[11pt]{article}
\usepackage[margin=1in]{geometry}
\usepackage{amsmath,amssymb}
\usepackage[hidelinks]{hyperref}

\title{Literature Status: the CAR Foguel--Hankel $B_2$ Condition}
\author{Run \texttt{fa\_banach\_001}, lane 3}
\date{19 July 2026}

\begin{document}
\maketitle

\paragraph{Status.}
\texttt{literature\_already\_answered}.

\paragraph{Original question.}
For a scalar sequence $\alpha=(\alpha_k)_{k\geq0}$, Badea's
\emph{Operators near completely polynomially dominated ones and similarity
problems} \cite{Badea} considers the associated CAR-valued Foguel--Hankel
operator $R(Y_\alpha)$ and sets
\[
 B_2(\alpha)=\sum_{k\geq0}(k+1)^2|\alpha_k|^2.
\]
Theorem 1.4 proves that similarity of $R(Y_\alpha)$ to a contraction implies
$B_2(\alpha)<\infty$.  Immediately afterward, on page 5 of the 22-page arXiv
PDF, the source asks:
\begin{quote}
It is an open problem if $B_2$ finite implies $R(Y_\alpha)$ similar to a
contraction.
\end{quote}
The question is repeated after Corollary 5.6 on page 14.

\paragraph{Identification of the answer.}
Ricard's distinct paper \emph{On a Question of Davidson and Paulsen}
\cite{Ricard} supplies the affirmative converse.  Its abstract explicitly
identifies the result as an answer to the Davidson--Paulsen question about
similarity to a contraction for this class of CAR-valued Foguel--Hankel
operators.  Thus, together with the already-known necessity,
\[
 R(Y_\alpha)\text{ is similar to a contraction}
 \quad\Longleftrightarrow\quad
 \sum_{k\geq0}(k+1)^2|\alpha_k|^2<\infty.
\]

A later open-access source gives independent full-text corroboration:
Chalmoukis--Stylogiannis \cite{CS}, Theorem 4.1 and the paragraph immediately
following it (PDF pages 14--15), state Ricard's Hankel Schur-multiplier theorem
and explicitly say that this theorem answers the Davidson--Paulsen question
about CAR-valued Foguel--Hankel operators similar to a contraction.

\paragraph{Verification note.}
Ricard's HAL record contains the DOI and journal metadata but currently no
deposited file, and the publisher full-text API requires a private key.  This
packet therefore does not claim direct local inspection of the journal PDF;
it relies on the answer-identifying journal abstract plus the directly
inspected later full-text attribution.  This is an explicit later-literature
answer, not a new theorem of this run.

\begin{thebibliography}{9}
\bibitem{Badea}
C. Badea,
\emph{Operators near completely polynomially dominated ones and similarity
problems}, Journal of Operator Theory \textbf{49} (2003), 3--23;
\href{https://arxiv.org/abs/math/0102160}{arXiv:math/0102160}.

\bibitem{Ricard}
E. Ricard,
\emph{On a Question of Davidson and Paulsen},
Journal of Functional Analysis \textbf{192} (2002), 283--294;
\href{https://doi.org/10.1006/jfan.2001.3899}{doi:10.1006/jfan.2001.3899}.

\bibitem{CS}
N. Chalmoukis and G. Stylogiannis,
\emph{Quasi-nilpotency of Generalized Volterra Operators on Sequence Spaces},
Results in Mathematics \textbf{76} (2021), 173;
\href{https://arxiv.org/abs/2005.01660}{arXiv:2005.01660}.
\end{thebibliography}

\end{document}

